% Source: FYDP_Summer/1.intro.tex (adapted to the journal format)
\section{Introduction}

Large language models can answer scientific questions at a level that makes them useful for tutoring and decision support. A model can still be confidently wrong, and it can state a wrong answer in fluent language that hides the error. Multi-agent debate reduces this risk because three agents answer the same question and review each other's reasoning over rounds~\cite{du2024improving}. Chain-of-thought prompting lets each agent expose the steps behind its answer, so peers can inspect the reasoning~\cite{wei2022chain}. The weak point sits at the end of the debate. Most systems combine the answers by majority vote, where every agent counts the same and the most frequent answer wins~\cite{wang-etal-2024-rethinking-bounds,choi2025debate}. A confident and wrong majority can then push a correct minority agent to drop its answer, and the vote records the wrong result as agreement. This failure is called sycophantic consensus~\cite{pitre-etal-2025-consensagent,he2026minority}. No current debate system checks an agent's claims against outside scientific evidence before that agent earns influence. This study builds a trust-calibrated multi-agent deliberation framework to close that gap. Each agent's answer is split into small factual claims. The claims are checked against retrieved scientific literature, and every agent holds a bounded trust score that rises with supporting evidence and falls with contradicting evidence. The final answer comes from trust-weighted aggregation instead of a head count. The method is tested on scientific reasoning datasets such as GPQA and MMLU-Pro, and it needs no fine-tuning because all models run as served inference.

Du et al. show that debate improves factuality and reasoning over a single pass, but their agents are treated equally and their claims are never checked against external sources~\cite{du2024improving}. Pitre et al. report that the correct answer is present in the discussion yet ignored in more than 20 percent of wrong-answer cases, and their prompt-rewriting fix still leaves the final decision on self-reported confidence~\cite{pitre-etal-2025-consensagent}. The uncertainty route of Yoffe et al. reweights agent influence during the debate, and their own ground-truth experiment shows the gain is capped by the quality of that self-reported signal~\cite{yoffe-etal-2025-debunc}. Token cost is the target of Fan et al., who decide when a debate is worth running, yet majority voting remains once the debate starts~\cite{fan2026imad}. Wang et al. combine many models through static aggregation and score well on open benchmarks, with no per-agent trust and no evidence check~\cite{wang2025mixture}. He et al. detect a suppressed correct minority after the debate and overturn the majority, which comes too late to protect the minority while it argues~\cite{he2026minority}.

Across these systems, influence is set by a vote, a self-reported score, a prompt rewrite, or a post-hoc flip, and none of them verifies an agent's claims against evidence while the debate is running. Self-reported confidence is a weak basis because training that makes models agreeable also makes them more sycophantic and less accurate~\cite{ibrahim2026training}, and alignment training has been shown to increase agreement with wrong user positions~\cite{chen2025self}. Theory also predicts that correlated agents drift to a wrong majority unless an outside correction enters the loop~\cite{estornell2024multillm}, and multi-agent systems fail in recurring ways that diagnosis alone does not fix~\cite{cemri2025why}. This study resolves the unsolved part by tying each agent's influence to external evidence during the debate. Claims are checked against retrieved passages from scientific sources~\cite{karpukhin2020dense,kushwaha2026agentbased}, a step in line with work that grounds correction in external tools~\cite{gou2024critic}. The trust score is bounded, so no agent is silenced or made all-powerful, and the final decision uses those weights. Scientific question sets are a suitable testbed because they carry verifiable answers and published literature that retrieval can reach, and GPQA and MMLU-Pro allow comparison against established baselines. The expected outcome is a reproducible framework that lowers consensus collapse, preserves a correct minority more often, keeps answer accuracy stable, and reports a calibrated trust signal.
